\section{Discussion and Limitations}

\subsection{Discussion}

The clearest lesson is that planner-to-executor contract design is an
architectural design variable, not a presentation detail. Under the
under-specified direct-action contract, the
planner often produces task-level verbs that are semantically plausible but not
dispatchable under the executor's tool namespace. Explicit contracts instead
declare the callable interface up front and reduce the need for executor-side
interpretation.

Runtime validation plays a complementary role. It does not replace a
well-specified contract, but it can intercept a bad first emission, surface a
concrete error, and recover the run with minimal additional machinery. That
pattern matters for embodied execution stacks that preserve
planner-to-executor separation, including middleware-style action interfaces and
guarded execution policies.

At a higher level, the studied boundary maps onto execution architectures
already familiar in robotics. The contract specifies the form of a
dispatchable request, analogous to how ROS~2 actions define typed request
fields and execution interfaces
\cite{biggs2019ros2actions,macenski2022ros2,martin2021plansys2}. Runtime
validation then decides whether a malformed or mismatched request should be
rejected before execution. A similar separation appears in behavior-tree and
TAMP systems, where precondition or guard checks determine whether a proposed
action can continue to the executor interface
\cite{colledanchise2018behaviortrees,behaviortreecppdocs,behaviortreecppconditions,garrett2020pddlstream,garrett2021integratedtamp}.

The P2 pattern suggests a second, narrower lesson. In the evaluated
comparisons, a
brief planner-side self-check leaves the success profile unchanged while
showing fewer planner/tool calls. The bounded interpretation is not that
self-checks are universally best, but that contract design and runtime design
remain separable design choices once the callable interface is explicit.

\subsection{Limitations}

The conclusions are bounded to a controlled simulator study. The majority of
the fixed evaluation matrix runs on a lightweight deterministic reference
backend rather than inside a live Isaac Sim physics stage; Isaac Sim serves as a limited qualitative cross-backend check. The
evidence therefore covers two task families and is limited to the current easy
and harder slices rather than a broad execution benchmark.

The action-interface ablation and runtime-only ablation are compact mechanism
checks rather than broad sweeps, and backend coverage differs across slices.
The runtime-only ablation is especially specific: it is designed around
recoverable invalid-first-action probes, so its strongest claim is about
lightweight repair at that boundary.

Using a deterministic mock planner means invalid-action patterns are
pre-specified boundary probes rather than stochastic errors from a real
foundation model. The empirical frame is therefore descriptive rather than
inferential; the deterministic backend and fixed task set preclude confidence
intervals or significance tests.

These limits still allow a concrete systems conclusion. In the tested setting,
planner-to-executor contracts and lightweight runtime safeguards materially
change invalid actions, recovery, and planner/tool call counts.
A natural next step is to replicate this factorial with stochastic LLM
planner backends and wider scene coverage to test whether the observed
ordering generalizes beyond injected boundary probes.
